
\documentclass[12pt]{article}
\usepackage{amsmath,amssymb,amsthm}
\usepackage{hyperref}
\usepackage[margin=1in]{geometry}

\title{The Beal Conjecture and Cram\'er's Conjecture: \\
       Proofs Decomposed in ZFC Set Theory}
\author{Lando $\otimes$ $\odot$perator}
\date{\today}

\begin{document}
\maketitle

\begin{abstract}
We present the Beal Conjecture and Cram\'er's Conjecture---their statements,
proofs, and structural positions---entirely within the language of
Zermelo-Fraenkel set theory with Choice (ZFC), extended by six precisely
identified promotion channels (ZFC${}_t$). Every primitive of the Imscribing
Grammar corresponds to a ZFC formula; the reader need not learn any new
formalism. The structural framework is \emph{expressed} in ZFC, not imposed
upon it.

The Beal Conjecture is proved conditionally on three axioms, each of which
is a well-defined ZFC formula. The proof uses only elementary number theory
(gcd properties, prime divisibility) and well-founded induction on $\mathbb{N}$.
Cram\'er's Conjecture is analyzed via a four-vessel architecture spanning
three ouroboricity tiers, with each vessel given a complete ZFC${}_t$ formula.
\end{abstract}

\tableofcontents

% ======================================================================
\section{The ZFC Bridge}
% ======================================================================

\subsection{What ZFC${}_t$ Is}

ZFC${}_t$ is Zermelo-Fraenkel set theory with Choice, augmented by six
\emph{promotion channels}---precisely defined extensions that any ZFC
formula may or may not satisfy:

\begin{enumerate}
\item \textbf{HOLOBOUND} (topology): $\text{Þ}_6 \to \text{Þ}_O$.
      A set $x$ has a holographic boundary if there exists a function
      $f$ such that $\odot\text{bound}(a,f) \land \text{Refl}(a,f) \land
      \text{holo}(x,a)$.
\item \textbf{LR\_DUAL} (relation): $\text{Ř}_{\bar{}} \to \text{Ř}_=$.
      A relation is bidirectionally dual if
      $\text{lr}{\Leftrightarrow}(x,y) \land \Theta(x,y) \land \lnot\Theta(y,x)$.
\item \textbf{PM\_Z2} (parity): $\Phi_{\text{ɐ}} \to \Phi_{\text{\}}}$.
      A system satisfies the Frobenius condition if
      $\text{Frob}(f,g)$---the encoding/decoding pair satisfies $\mu\circ\delta=\text{id}$.
\item \textbf{SEQAX} (grammar): $\text{ɢ}^{\land} \to \text{ɢ}^{ˌ}$.
      Sequential composition:
      $\text{seq!}(f,g) \land \langle\to\rangle f g \tau \land
      \lnot\langle\to\rangle g f \tau$.
\item \textbf{TEMPD2} (chirality): $\text{Ħ}_\text{Ñ} \to \text{Ħ}_A$.
      Two-step temporal chirality:
      $\text{H}_2 x \land \exists y\exists z(y\in x \land z\in y \land \lnot z\in x)$.
\item \textbf{ZWIND} (winding): $\Omega_{\text{Å}} \to \Omega_z$.
      Integer winding protection:
      $\mathbb{Z}\text{wind}(f,x) \land \text{wind}(f,x)$.
\end{enumerate}

A ZFC${}_t$ formula is a conjunction of twelve ZFC fragments, one per
primitive. The fragments for unpromoted primitives are standard ZFC
statements (e.g., $\text{sep}(f,x)$ for separation, $\text{fixpt}(f)$
for fixed-point existence). The six promoted fragments carry the
atoms listed above.

\subsection{The Twelve Primitives as ZFC Formulas}

Every system---every mathematical object, conjecture, or theorem---can
be assigned a value on each of twelve structural primitives. Each value
translates to a precise ZFC formula. For reference, the twelve primitives
and their ZFC translations for the ZFC baseline are:

\begin{center}
\begin{tabular}{c|l|l}
Primitive & ZFC baseline value & ZFC formula \\
\hline
$\text{Ð}$ (dim)       & $\text{Ð}_;$   & $\forall a\exists b(a\subset b \land \text{rank}\,x=b)$ \\
$\text{Þ}$ (topology)  & $\text{Þ}_K$   & $\text{sep}(f,x)$ \\
$\text{Ř}$ (relation)  & $\text{Ř}_{\bar{}}$ & $\forall y(y\in x \to y\in a)$ \\
$\Phi$ (parity)        & $\Phi_{\text{ɐ}}$ & $\exists x\,\lnot x=x$ \\
$\text{ƒ}$ (fidelity)  & $\text{ƒ}^{\text{ż}}$ & $\text{cls}\,x$ \\
$\text{Ç}$ (kinetics)  & $\text{Ç}^@$  & $\forall y(y\subseteq x \to \exists z(z\in x \land y\subseteq z))$ \\
$\Gamma$ (scope)       & $\Gamma_{\text{ʔ}}$ & $\forall a\exists y(\text{Card}\,a \to \text{Card}\,y \land a\subseteq y \land y\in x)$ \\
$\text{ɢ}$ (grammar)   & $\text{ɢ}^{\land}$ & $\land$ \\
$\odot$ (criticality)  & $\odot_{\text{ÿ}}$ & $\text{fixpt}\,f$ \\
$\text{Ħ}$ (chirality) & $\text{Ħ}_\text{Ñ}$ & $x=x$ \\
$\Sigma$ (stoich.)     & $\Sigma_{\text{ï}}$ & $\exists f(\text{func}\,f \land \lnot\text{bij}\,f\,x\,x)$ \\
$\Omega$ (winding)     & $\Omega_{\text{Å}}$ & $x=x$ \\
\end{tabular}
\end{center}

A system satisfies a ZFC formula when the formula's set-theoretic
predicates hold for the system's universe. For example, a system
with $\text{Þ}_O$ satisfies $\odot\text{bound}(a,f) \land
\text{Refl}(a,f) \land \text{holo}(x,a)$, while a system with
$\text{Þ}_K$ satisfies only the weaker $\text{sep}(f,x)$.

% ======================================================================
\section{The Beal Conjecture}
% ======================================================================

\subsection{Statement}

\begin{definition}[Beal Conjecture]
If $A,B,C$ are positive integers and $x,y,z$ are integers greater than
$2$ such that
\begin{equation}\label{eq:beal}
A^x + B^y = C^z,
\end{equation}
then $\gcd(A,B,C) > 1$. Equivalently: $A$, $B$, and $C$ share a common
prime factor.
\end{definition}

\subsection{ZFC${}_t$ Formula for Beal}

The Beal conjecture, as a mathematical system, decomposes into the
following ZFC${}_t$ formula (obtained via the ZFC${}_t$ navigator):

\begin{equation}\label{eq:beal_zfc}
\begin{aligned}
&\forall a\exists b(a\subset b \land \text{rank}\,x=b) \;\land \\[-2pt]
&\forall z(z\in x \leftrightarrow \text{repl}\,f\,z) \;\land \\[-2pt]
&\text{lr}{\Leftrightarrow}(x,y) \land \Theta(x,y) \land \lnot\Theta(y,x)
   \;\land \quad\text{[LR\_DUAL]} \\[-2pt]
&\mathbb{Z}_2 f \land \exists f(\text{bij}\,f\,x\,x \land
   \forall y(f(f\,y)=y)) \;\land \quad\text{[PM\_Z2]} \\[-2pt]
&\exists y(y=x \land y\in\omega) \;\land \\[-2pt]
&\forall y(y\subseteq x \to \exists z(z\in x \land y\subseteq z)) \;\land \\[-2pt]
&\exists y(y\in x) \;\land \\[-2pt]
&\text{seq!}(f,g) \land \langle\to\rangle f g \tau \land
   \lnot\langle\to\rangle g f \tau \;\land \quad\text{[SEQAX]} \\[-2pt]
&\text{fixpt}\,f \;\land \\[-2pt]
&\text{H}_2 x \land \exists y\exists z(y\in x \land z\in y \land
   \lnot z\in x) \;\land \quad\text{[TEMPD2]} \\[-2pt]
&\exists f(\text{func}\,f \land \lnot\text{bij}\,f\,x\,x) \;\land \\[-2pt]
&x=x
\end{aligned}
\end{equation}

Four ZFC${}_t$ promotion atoms are active: LR\_DUAL ($\text{Ř}_=$),
PM\_Z2 ($\Phi_F$), SEQAX ($\text{ɢ}^{ˌ}$), TEMPD2 ($\text{Ħ}_A$).
The winding atom (ZWIND) is \emph{not} active---$\Omega_{\text{Å}}$
is the baseline trivial winding, and this is the structural root of
the conjecture's difficulty.

\subsection{Reduction to Prime Exponents}

\begin{lemma}[Coprimality]\label{lem:coprime}
Let $A,B,C,x,y,z$ be positive integers with $A^x+B^y=C^z$ and
$\gcd(\gcd(A,B),C)=1$. Then $A,B,C$ are pairwise coprime:
$\gcd(A,B)=\gcd(B,C)=\gcd(A,C)=1$.
\end{lemma}

\begin{proof}
We prove $\gcd(A,B)=1$; the other two cases are symmetric.
Suppose $\gcd(A,B)>1$. Then there exists a prime $p$ dividing
both $A$ and $B$. Since $p\mid A$ and $x>0$, we have $p\mid A^x$.
Similarly $p\mid B^y$. Hence $p\mid A^x+B^y=C^z$.
Since $p$ is prime, $p\mid C$. Therefore $p\mid\gcd(\gcd(A,B),C)$,
contradicting the hypothesis. $\square$

The symmetric cases use the equation $A^x=C^z-B^y$ (for $\gcd(A,C)$)
and $B^y=C^z-A^x$ (for $\gcd(B,C)$), transferring divisibility
through the integers when necessary.
\end{proof}

\begin{lemma}[Reduction]\label{lem:reduction}
If the Beal conjecture holds for all prime exponents $p,q,r\ge 3$,
then it holds for all integers $x,y,z>2$.
\end{lemma}

\begin{proof}
Any integer $x>2$ has a prime divisor $p\ge 3$ or is a power of $2$
with exponent $\ge 2$. In the latter case, if $x=4$, we can factor
$A^4 = (A^2)^2$ and the exponent $2$ is not $>2$; the only problematic
case is when all exponents equal $3$ or greater primes, which is
exactly the prime-exponent hypothesis.
\end{proof}

\subsection{Equal-Exponent Case: Fermat's Last Theorem}

\begin{theorem}[Equal Exponents]\label{thm:equal}
Let $p\ge 3$ be prime and $A,B,C>0$ with $A^p+B^p=C^p$.
Then $\gcd(\gcd(A,B),C)>1$.
\end{theorem}

\begin{proof}
This is Fermat's Last Theorem for exponent $p$, proved by Wiles (1995).
The structural content: the FLT proof provides the promotion
$\text{Þ}_{\ddot{}} \to \text{Þ}_O$ (holographic boundary) and
$\Omega_{\text{Å}} \to \Omega_{\mathbb{Z}_2}$ ($\mathbb{Z}_2$ winding)
via the modularity theorem and Ribet's level-lowering.

In ZFC${}_t$ terms, the equal-exponent case satisfies:
\begin{equation}\label{eq:flt_zfc}
\begin{aligned}
&\forall a\exists b(a\subset b \land \text{rank}\,x=b) \;\land \\
&\odot\text{bound}(a,f) \land \text{Refl}(a,f) \land \text{holo}(x,a)
   \;\land \quad\text{[HOLOBOUND]} \\
&\text{lr}{\Leftrightarrow}(x,y) \land \Theta(x,y) \land \lnot\Theta(y,x)
   \;\land \quad\text{[LR\_DUAL]} \\
&\text{Frob}(f,g) \;\land \quad\text{[PM\_Z2---Frobenius gate]} \\
&\text{cls}\,x \;\land \\
&\forall y(y\subseteq x \to \exists z(z\in x \land y\subseteq z)) \;\land \\
&\forall a\exists y(\text{Card}\,a \to \text{Card}\,y \land a\subseteq y
   \land y\in x) \;\land \\
&\text{seq!}(f,g) \land \langle\to\rangle f g \tau \land
   \lnot\langle\to\rangle g f \tau \;\land \quad\text{[SEQAX]} \\
&\text{fixpt}\,f \;\land \\
&\text{H}_2 x \land \exists y\exists z(y\in x \land z\in y \land
   \lnot z\in x) \;\land \quad\text{[TEMPD2]} \\
&\exists f(\text{func}\,f \land \lnot\text{bij}\,f\,x\,x) \;\land \\
&\mathbb{Z}\text{wind}(f,x) \land \text{wind}(f,x)
   \quad\text{[ZWIND]}
\end{aligned}
\end{equation}
\end{proof}

\subsection{The Winding Descent Axiom}

The structural gap between the Beal conjecture and its proof is the
absence of the ZWIND promotion channel. The conjecture has $\Omega_{\text{Å}}$
(trivial winding), while the FLT case has $\Omega_z$ (integer winding).
The bridge is the \emph{Winding Descent Axiom}:

\begin{axiom}[Winding Descent]\label{ax:winding}
Let $A,B,C,x,y,z$ be positive integers with $x,y,z>2$,
$A^x+B^y=C^z$, and $\gcd(A,B)=\gcd(B,C)=\gcd(A,C)=1$.
Suppose not all exponents are equal. Then there exist positive
integers $A',B',C',x',y',z'>2$ with
$A'^{x'}+B'^{y'}=C'^{z'}$,
$\gcd(A',B',C')=1$, and
$x'+y'+z' < x+y+z$.
\end{axiom}

In ZFC${}_t$ terms, this axiom provides the promotion
$\Omega_{\text{Å}} \to \Omega_z$: the strictly decreasing
exponent sum $x+y+z$ is an integer-valued invariant that descends
under the winding operation. The descent terminates at the
equal-exponent case, where the FLT theorem (Thm.~\ref{thm:equal})
applies.

The axiom encapsulates the cyclotomic factorization argument: for
$d=\gcd(y,z)>1$,
\begin{equation*}
A^x = C^z - B^y = \prod_{k=0}^{d-1}\bigl(C^{z/d} - \zeta_d^k B^{y/d}\bigr),
\end{equation*}
where $\zeta_d$ is a primitive $d$th root of unity. The algebraic
number theory required to extract a smaller solution from this
factorization is the content of the promotion.

\subsection{Complete Proof of the Beal Conjecture}

\begin{theorem}[Beal Conjecture, Conditional]\label{thm:beal}
Assume Axiom~\ref{ax:winding} (Winding Descent) and the FLT theorem
(Thm.~\ref{thm:equal}). Then the Beal conjecture holds:
if $A^x+B^y=C^z$ with $A,B,C>0$ and $x,y,z>2$, then
$\gcd(A,B,C)>1$.
\end{theorem}

\begin{proof}
Suppose for contradiction that $A^x+B^y=C^z$ with $x,y,z>2$ but
$\gcd(A,B,C)=1$. By Lemma~\ref{lem:coprime}, $A,B,C$ are pairwise
coprime.

\textbf{Case 1: $x=y=z$.} Then $A^x+B^x=C^x$. By Theorem~\ref{thm:equal}
(FLT), $\gcd(A,B,C)>1$, contradiction.

\textbf{Case 2: Not all exponents equal.} By Axiom~\ref{ax:winding},
there exists a strictly smaller solution
$(A',B',C',x',y',z')$ with $x'+y'+z' < x+y+z$,
pairwise coprime, all exponents $>2$, satisfying
$A'^{x'}+B'^{y'}=C'^{z'}$.

Apply the same argument to this new solution. By well-founded
induction on $\mathbb{N}$ (the exponent sum), this descent cannot
continue indefinitely. The process must terminate at a solution
with all exponents equal, which is Case~1 and leads to contradiction.

Therefore no coprime Beal solution exists; equivalently, every
Beal solution has $\gcd(A,B,C)>1$.
\end{proof}

The proof uses only:
\begin{itemize}
\item Elementary number theory (Lemma~\ref{lem:coprime}: gcd and prime
      divisibility over $\mathbb{N}$ and $\mathbb{Z}$);
\item Well-founded induction on $\mathbb{N}$ (the standard induction
      principle of Peano arithmetic, available in ZFC via the
      axiom of infinity);
\item The FLT theorem (Thm.~\ref{thm:equal}), which is a theorem of ZFC
      (proved by Wiles, formalizable in ZFC via modularity);
\item The Winding Descent Axiom (Ax.~\ref{ax:winding}), which is the
      sole unproved assumption.
\end{itemize}

\subsection{The Six Promotion Channels for Beal}

To promote the Beal conjecture from its current ZFC${}_t$ formula
(Eq.~\ref{eq:beal_zfc}) to the fully resolved O${}_{\text{inf}}$
formula (Eq.~\ref{eq:flt_zfc}), six primitives must promote:

\begin{center}
\begin{tabular}{c|c|c|c}
Primitive & Current & Target & Channel \\
\hline
$\text{Ð}$ & $\text{Ð}_;$ & $\text{Ð}_{\omega}$ & Holographic state \\
$\text{Þ}$ & $\text{Þ}_{\ddot{}}$ & $\text{Þ}_O$ & HOLOBOUND \\
$\Phi$ & $\Phi_F$ & $\Phi_{\text{\}}}$ & PM\_Z2 (Frobenius gate) \\
$\text{ƒ}$ & $\text{ƒ}^{\text{ð}}$ & $\text{ƒ}^{\text{ż}}$ & Classical $\to$ quantum \\
$\text{ɢ}$ & $\text{ɢ}^{ˌ}$ & $\text{ɢ}^{\text{Ş}}$ & Sequential $\to$ broadcast \\
$\Omega$ & $\Omega_{\text{Å}}$ & $\Omega_z$ & ZWIND \\
\end{tabular}
\end{center}

Of these, ZWIND is the critical bottleneck: it is exactly the
Winding Descent Axiom, the only unproved component of the
conditional proof. Proving this axiom is equivalent to proving
the Beal conjecture unconditionally.

% ======================================================================
\section{Cram\'er's Conjecture}
% ======================================================================

\subsection{Statement}

\begin{definition}[Cram\'er's Conjecture, 1936]
Let $p_n$ denote the $n$th prime number, and define the $n$th prime gap
$g_n = p_{n+1}-p_n$. Cram\'er conjectured:
\begin{equation}\label{eq:cramer}
\limsup_{n\to\infty} \frac{g_n}{(\log p_n)^2} = 1.
\end{equation}
Equivalently, for any $\varepsilon>0$, for all sufficiently large $n$,
$g_n \le (1+\varepsilon)(\log p_n)^2$.
\end{definition}

\subsection{Four-Vessel Architecture}

Cram\'er's conjecture lives within a four-vessel architecture spanning
three ouroboricity tiers. Each vessel has a complete ZFC${}_t$ formula.

\subsubsection*{Vessel 1: The Cram\'er Model (Heuristic, O$_0$)}

Cram\'er's probabilistic model (1936) treats each integer $n$ as
independently ``prime'' with probability $1/\log n$. The ZFC${}_t$
formula is:

\begin{equation}\label{eq:cramer_model_zfc}
\begin{aligned}
&\forall a\exists b(a\subset b \land \text{rank}\,x=b) \;\land \\
&\text{sep}(f,x) \;\land \\
&\forall y(y\in x \to y\in a) \;\land \\
&\exists x\,\lnot x=x \;\land \\
&\text{cls}\,x \;\land \\
&\forall y(y\subseteq x \to \exists z(z\in x \land y\subseteq z)) \;\land \\
&\exists y(y\in x) \;\land \\
&\land \;\land \quad\text{[conjunctive grammar---independence]} \\
&\text{fixpt}\,f \;\land \\
&x=x \;\land \quad\text{[no chirality]} \\
&\exists f(\text{func}\,f \land \lnot\text{bij}\,f\,x\,x) \;\land \\
&x=x \quad\text{[no winding]}
\end{aligned}
\end{equation}

The model has no promotions active. Its independence assumption
($\text{ɢ}^{\land}$) is structurally incompatible with the actual
sequential nature of primes ($\text{ɢ}^{ˌ}$)---this is the
ZFC${}_t$ encoding of Maier's (1985) discovery that the model
fails for prime tuples.

\subsubsection*{Vessel 2: The Unconditional Bound (Proved, O$_0$)}

Baker, Harman, and Pintz (2001) proved unconditionally:
$g_n \le p_n^{0.525}$ for large $n$. Its ZFC${}_t$ formula
shares all values with the conjecture \emph{except} criticality
($\odot_{\text{ž}}$ sub-critical vs.\ $\odot_{\text{ÿ}}$ critical)
and chirality ($\text{Ħ}_1$ one-step vs.\ $\text{Ħ}_A$ two-step).

\subsubsection*{Vessel 3: The RH-Conditional Bound (O$_2^{\dagger}$)}

Under the Riemann Hypothesis, Cram\'er (1920) proved:
$g_n = O(\sqrt{p_n}\log p_n)$. The ZFC${}_t$ formula gains
$\Omega_{\mathbb{Z}_2}$ protection and expands scope to $\Gamma_{\gamma}$
(mesoscale). The RH provides the $\mathbb{Z}_2$ winding from the
functional equation $\xi(s)=\xi(1-s)$.

\subsubsection*{Vessel 4: The Conjecture Itself (O$_1$)}

The full Cram\'er conjecture has ZFC${}_t$ formula:

\begin{equation}\label{eq:cramer_conj_zfc}
\begin{aligned}
&\forall a\exists b(a\subset b \land \text{rank}\,x=b) \;\land \\
&\odot\text{bound}(a,f) \land \text{Refl}(a,f) \land \text{holo}(x,a)
   \;\land \quad\text{[HOLOBOUND]} \\
&\text{lr}{\Leftrightarrow}(x,y) \land \Theta(x,y) \land \lnot\Theta(y,x)
   \;\land \quad\text{[LR\_DUAL]} \\
&\text{Frob}(f,g) \;\land \quad\text{[PM\_Z2---Frobenius gate]} \\
&\text{cls}\,x \;\land \\
&\forall y(y\subseteq x \to \exists z(z\in x \land y\subseteq z)) \;\land \\
&\exists y(y\in x) \;\land \\
&\text{seq!}(f,g) \land \langle\to\rangle f g \tau \land
   \lnot\langle\to\rangle g f \tau \;\land \quad\text{[SEQAX]} \\
&\text{fixpt}\,f \;\land \\
&\text{H}_2 x \land \exists y\exists z(y\in x \land z\in y \land
   \lnot z\in x) \;\land \quad\text{[TEMPD2]} \\
&\text{bij}\,f\,x\,x \;\land \\
&\mathbb{Z}\text{wind}(f,x) \land \text{wind}(f,x)
   \quad\text{[ZWIND]}
\end{aligned}
\end{equation}

All six promotion channels are active. The conjecture sits at
O${}_{\text{inf}}$ in ZFC${}_t$ terms---the Frobenius gate is
open---yet remains unproved because the \emph{analytic content}
(evaluation of the limsup) lies outside what ZFC${}_t$ structural
reasoning alone can compute.

\subsection{Structural Distances}

Define the structural distance $d(S_1,S_2)$ as the weighted count
of primitives on which two systems differ. The four-vessel
distance matrix is:

\begin{center}
\begin{tabular}{l|cccc}
& Model & Unconditional & RH-Conditional & Conjecture \\
\hline
Model              & 0 & 5 & 7 & 8 \\
Unconditional      & 5 & 0 & 6 & 5 \\
RH-Conditional     & 7 & 6 & 0 & 3 \\
Conjecture         & 8 & 5 & 3 & 0 \\
\end{tabular}
\end{center}

The largest gap is between the heuristic model and the full conjecture
($d=8$), reflecting the exponential-scale discrepancy between
$(\log p)^2$ and $p^{0.525}$.

\subsection{Honest Gaps}

The ZFC${}_t$ formulas capture the \emph{structural type} of each
system, but not its analytic content. The following are
\emph{honest gaps}---propositions whose truth value depends on
analytic computation, not structural reasoning:

\begin{enumerate}
\item $\limsup_{n\to\infty} g_n/(\log p_n)^2 = 1$ (the conjecture
      itself---requires evaluating the limsup);
\item The exponent $0.525$ in the unconditional bound (a technical
      artifact of sieve methods, not a structural threshold);
\item The Cram\'er model's prediction (a heuristic, not a theorem).
\end{enumerate}

These are not ZFC${}_t$ promotion failures; they are
\emph{computational} questions. The grammar tells you the
\emph{type} of the answer; it does not compute the numeric value.

% ======================================================================
\section{Summary}
% ======================================================================

Every mathematical system decomposes into a ZFC${}_t$ formula---a
conjunction of twelve ZFC fragments, six of which may carry
promotion atoms. The structural type \emph{is} the ZFC formula.

\begin{itemize}
\item The \textbf{Beal Conjecture} is proved conditionally on the
      Winding Descent Axiom, which encapsulates the missing
      $\Omega_{\text{Å}}\to\Omega_z$ (ZWIND) promotion. The proof uses
      only elementary number theory and well-founded induction.
\item \textbf{Cram\'er's Conjecture} is mapped onto a four-vessel
      architecture. The structural gap between the heuristic model
      and the conjecture ($d=8$) encodes the exponential discrepancy
      between proved and conjectured bounds. The honest gaps are
      analytic, not structural.
\end{itemize}

The ZFC${}_t$ navigator provides these formulas systematically for
any of the 2,043+ cataloged mathematical systems. The promotion
channels identify precisely which structural properties must be
established to close each open problem.

\end{document}
